% R15 component; only input paths and enumerated location/grammar prose are edited.
% Source: revisions/2026-10-04-r15/manuscript/capital.tex; commit 1cb3cc9efe135966ec228dfaf51c6841a6dfc97c.
% Generated by code/integrate.py; edit extraction rules there.
% Reviewed source: revisions/2026-10-04-r10/manuscript/continuous.tex, line 1, commit f5021cefa71492babfcbfa580e0c984f59a9de26
\section{Continuous-Economy Guarantees for Neural Capital Policies}\label{sec:r10continuous}
The distinction between a nodal certificate and an economic diffusion bound is substantive. This section obtains the latter for the original dense, nonlinear capital economy, without replacing its primitives, restricting the comparison policy class, or extrapolating a sampled residual. The construction uses the curvature of consumption utility to control the gain from an arbitrary adapted deviation. A neural actor is then confined to a known neighbourhood of a feasible schedule. Its approximation class is restricted; the economic optimum against which it is assessed is not.

Write $\bar z=d^{-1}\sum_i z_i$, $\|z\|_d^2=d^{-1}\sum_i z_i^2$, and $\Pi z=z-\bar z\bm1$. In the capital model already specified, the state and payoff are
\begin{align}
 dY_i&=\{a+\kappa\tanh((BY)_i)-m_i-\nu^2/2\}\,dt
       +\sigma_I\,dW_i+\sigma_C\,dW_0,\label{eq:r10state}\\
 J_y(m)&=\E_y\left[\int_0^T e^{-\rho t}
  \{\overline{\log m_t}+\bar Y_t-\tfrac\zeta2\bar m_t^2\}\,dt
  +e^{-\rho T}\{\bar Y_T-\chi\|\Pi Y_T\|_d^2\}\right],\label{eq:r10payoff}
\end{align}
where $\nu^2=\sigma_I^2+\sigma_C^2$ and each control lies in $[\underline m,M]$. Controls are progressively measurable and the state equation is interpreted on $\R^d$. Bounded control and the globally Lipschitz bounded production term yield the required state moments and finite payoffs. The historical parameters are $T=1$, $\rho=0.04$, $a=0.10$, $\kappa=0.15$, $\sigma_I=0.15$, $\sigma_C=0.10$, $\zeta=0.2$, $\chi=0.03$, and $[\underline m,M]=[0.02,2]$. The stored coupling matrices are unchanged in dimensions ten and twenty. Dimension fifty uses the same stated matrix-generation rule.

Define
\begin{align}
 w(t)&=\frac{1-e^{-\rho(T-t)}}{\rho}+e^{-\rho(T-t)},&
 m_0(t)&=\frac{2}{w(t)+\sqrt{w(t)^2+4\zeta}}.\label{eq:r10anchor}
\end{align}
This schedule is strictly interior for the parameters above. It is a feasible benchmark, not an asserted solution of the coupled economy. 
For the analytical comparison, keep $m_0$ interior and assume
$\rho>0$, $\zeta,\chi,\kappa\geq0$,
and let $\beta\geq\kappa\|B\|_2$ and $s_0\geq\|\Pi y\|_d$.
Choose $\epsilon\geq0$ so that
$[m_0(t)-\epsilon,m_0(t)+\epsilon]\subset[\underline m,M]$
at every date.
With $V^*(y)=\sup_m J_y(m)$ over the full adapted admissible class,
the schedule gap obeys
$0\leq V^*(y)-J_y(m_0)\leq G_d$. The exact scalar integral defining $G_d$
is \eqref{eq:r10anchorbound}; Appendix~\ref{app:r15capitalaccounts}
gives its curvature and dispersion constants, Theorem~\ref{thm:r10tube},
and the complete proof. This analytical bound supplies a common economic
reference. The policy-specific endpoint below measures what the fitted actor
adds to that reference.

% Reviewed source: revisions/2026-10-04-r10/manuscript/continuous.tex, line 53, commit f5021cefa71492babfcbfa580e0c984f59a9de26
\subsection{The deployed neural map and its numerical guard}
The actor is a two-hidden-layer network with the closed-form output map
\begin{equation}\label{eq:r10actor}
 m_i^\phi(t,Y)=m_0(t)+\epsilon\tanh N_{\phi,i}(t,Y).
\end{equation}
The weights are trained by the same block-specific NBO losses as before. No maximizing policy labels, nodal action table, or online local search enter this actor. Input clipping at very large log-capital values protects numerical evaluation; it does not stop or truncate the economic state. The theorem remains valid for any bounded proposal followed by the stated output constraint, including this clipped-input network.

Numerical implementation must not replace the exact schedule by an unverified floating-point centre. The deployment routine therefore encloses $m_0(t)$ by $[\ell(t),u(t)]$ using the inherited outward arithmetic kernel and projects the proposal onto
\begin{equation}\label{eq:r10inward}
 [\,\operatorname{up}(u(t)-\epsilon),\,
       \operatorname{down}(\ell(t)+\epsilon)\,].
\end{equation}
This interval lies inside the exact tube. Adjacent clock floats enclose rounding of the input time. An empty interval causes an error; a nonfinite proposal is replaced by its interval midpoint and is not evidence of successful training. This scalar feasibility guard is neither an action-maximization oracle nor a stored policy correction table. Its cost is included in the deployment measurements.

% Reviewed source: revisions/2026-10-04-r11/manuscript/policy_specific.tex, line 1, commit f5021cefa71492babfcbfa580e0c984f59a9de26
\subsection{Policy-Specific Verification after Bellman Improvement}\label{sec:r11specific}
The bound in Theorem~\ref{thm:r10tube} is a useful starting point, but not a measure of what training accomplishes. Every policy in the same tube receives the same allowance, including a poorly trained policy. We now replace that allowance by an independently verified, policy-specific payoff difference. The economic primitives, initial state, and full adapted comparison class in \eqref{eq:r10state}--\eqref{eq:r10payoff} are unchanged. The implementation to which the new bound applies is specified explicitly below. It is not identified with continuously observed-state feedback.

\subsubsection{The deployed controller}
Fix a deterministic grid $t_k=kh$, $Nh=T$. The controller has an internal state $Z_k$, initially $y$, and observes the Brownian innovations at the grid dates. At date $t_k$ it selects
\begin{equation}\label{eq:r11implementation}
 a_k=\mathcal P_k\{m_\phi(t_k,Z_k)\},\qquad
 Z_{k+1}=Z_k+h\{c\bm1+f(Z_k)-a_k\}+\Sigma\Delta\widetilde W_k,
\end{equation}
where $c=a-\nu^2/2$, $f(z)=\kappa\tanh(Bz)$, and $\mathcal P_k$ projects inward onto the declared action interval around $m_0(t_k)$. The vector $\Delta\widetilde W_k$ clips each standardized innovation at ten before rescaling by $\sqrt h$. The numerical implementation and its rounding convention form part of the controller. Between dates it holds $a_k$ fixed. Its economic state, denoted $Y^\phi$, nevertheless satisfies
\[
 dY^\phi_t=\{c\bm1+f(Y^\phi_t)-a_k\}\,dt+\Sigma\,dW_t,
 \qquad t_k\leq t<t_{k+1}.
\]
In particular, the economic Brownian motion and state are not clipped, stopped, or replaced by $Z$. The controller is adapted to the original innovation filtration and remains feasible. Observing innovations is an implementation assumption, not a conclusion inferred from a state-only observation model. The continuously evaluated absolute-bound study policies remain separate objects in the supplement.

This specification has an important analytical consequence. When the numerical state and the economic state are compared, they use exactly the same chosen actions. Their state discrepancy therefore depends on the smoothness of the economic drift, not on a global Lipschitz bound for the learned network. The following construction uses this fact rather than inserting sampled residuals into a uniform verification theorem.

\subsubsection{An exact improvement identity}
Write $W(t)=e^{-\rho t}w(t)$ and define the nonnegative consumption deficit
\begin{equation}\label{eq:r11deficit}
 R_t(m)=\log m_0(t)-\overline{\log m}
       +w(t)\{\bar m-m_0(t)\}
       +\frac\zeta2\{\bar m^2-m_0(t)^2\}.
\end{equation}
The first-order identity $m_0^{-1}=w+\zeta m_0$ and concavity of the logarithm imply $R_t(m)\geq0$. The exact mean-state identity in the proof of Theorem~\ref{thm:r10tube} gives
\begin{align}
 \Delta_\phi:=J_y(a^\phi)-J_y(m_0)
 ={}&\E\int_0^T W(t)\,\overline{f(Y^\phi_t)-f(Y^0_t)}\,dt
       -\E\int_0^T e^{-\rho t}R_t(a^\phi_t)\,dt\notag\\
 &-\chi e^{-\rho T}\E\{\|\Pi Y^\phi_T\|_d^2-\|\Pi Y^0_T\|_d^2\}.
 \label{eq:r11gain}
\end{align}
Thus state-dependent consumption has to earn its utility-curvature cost through production or terminal dispersion. The identity is also an effective variance reduction: the common mean-state martingale cancels before simulation statistics are formed. Positive neural output is not itself evidence of positive economic improvement.

For each cell, calculate the scalar integrals
\begin{align*}
 A_k&=\int_{t_k}^{t_{k+1}}e^{-\rho t}dt,&
 B_k&=\int_{t_k}^{t_{k+1}}W(t)dt,\\
 C_k&=\int_{t_k}^{t_{k+1}}e^{-\rho t}
       \{\log m_0(t)-w(t)m_0(t)-\zeta m_0(t)^2/2\}\,dt.
\end{align*}
The consumption-deficit integral for a held action is exactly
$C_k-A_k\overline{\log a_k}+B_k\bar a_k+\zeta A_k\bar a_k^2/2$.
Consequently, consumption need not be approximated by a left-point utility sum. Let $X_{\phi,h}$ be the paired statistic obtained from \eqref{eq:r11gain} by using the internal Euler states in the production and terminal terms and these integrated consumption terms. The reference internal state uses the integrated common action $\int_{t_k}^{t_{k+1}}m_0(t)dt$. All scalar integral enclosures and their midpoint errors are saved.


\subsubsection{Transferring the computed payoff to the economy}
The statistic $X_{\phi,h}$ is evaluated for the held controller in
\eqref{eq:r11implementation}, using the same actions when comparing its
internal state with its economic diffusion. Proposition~\ref{prop:r11transfer}
in Appendix~\ref{app:r15capitalaccounts} bounds
$|\Delta_\phi-\E X_{\phi,h}|$ by the saved allowance
$E_{\phi,h}^{\rm all}$. The exact-arithmetic diffusion allowance is
$E_\epsilon(h)+E_0(h)$, with $E_\epsilon$ defined in
\eqref{eq:r11bias}. The appendix gives the additional clipped-innovation,
forward-arithmetic, and enclosed-integration terms. The $O(h)$ diffusion
bound uses bounded actions, the global drift and derivative bounds of the
specified hyperbolic-tangent production function, and additive diffusion;
it requires no derivative bound for the trained network.

\subsubsection{A post-training certificate}
Use the deterministic clipping threshold $b_\phi$ and expectation-tail
allowance $\tau_\phi$ constructed in
Appendix~\ref{app:r15capitaltail}. They account for the unbounded Gaussian
terminal-dispersion term and are not estimated from the final sample range.
Let $X^c_{\phi,h}$ be $X_{\phi,h}$ clipped to
$[-b_\phi,b_\phi]$, so that
$|\E X_{\phi,h}-\E X^c_{\phi,h}|\leq\tau_\phi$.

% Reviewed source: revisions/2026-10-04-r11/manuscript/policy_specific.tex, line 98, commit f5021cefa71492babfcbfa580e0c984f59a9de26
For $n$ fresh independent paths, let $\bar X^c_\phi$ and $s^2_\phi$ be the clipped empirical mean and unbiased variance. Fix a family of at most $F$ one-sided statements before inspecting this noise bank, and set $\ell=\log(2F/\alpha)$. Define
\begin{equation}\label{eq:r11lower}
 L_\phi=\bar X^c_\phi-
     \sqrt{2s_\phi^2\ell/n}-\frac{14b_\phi\ell}{3(n-1)}
     -E_{\phi,h}^{\rm all}-\tau_\phi.
\end{equation}
In the implementation the empirical mean is enclosed downward and the variance and every subtracted quantity upward. The upper payoff endpoint uses the same terms with reversed signs. The factor $F$ counts both tails when both are reported.

\Needspace{8\baselineskip}
\begin{theorem}[Policy-specific continuous-time certificate]\label{thm:r11certificate}
Condition on the training and validation data, and suppose the specified final simulation paths are independent draws from the declared innovation law. Under the arithmetic contract and transfer assumptions above, with probability at least $1-\alpha$, simultaneously for all declared policies and initial states,
\begin{equation}\label{eq:r11regret}
 J_y(a^\phi)-J_y(m_0)\geq L_\phi,
 \qquad 0\leq V^*(y)-J_y(a^\phi)\leq G_d-L_\phi.
\end{equation}
The optimum $V^*$ still ranges over all adapted feasible controls in the original continuous economy, not merely the networks, schedules, or sampled controllers used in the experiment. Common shocks may be reused across policies; independence across those policies is not required. The guarantee is pointwise in each declared initial state and conditional on training, rather than uniform on an uncountable state domain.
\end{theorem}

\begin{proof}
Apply the empirical Bernstein inequality of \citet[Theorem~4]{maurer2009} to $(X^c_{\phi,h}+b_\phi)/(2b_\phi)$ and use a union bound over the declared one-sided statements. Rescaling gives the first two error terms in \eqref{eq:r11lower}. The clipping expectation bound and Proposition~\ref{prop:r11transfer} then give the first inequality in \eqref{eq:r11regret}. Subtract it from the deterministic anchor bound \eqref{eq:r10anchorbound}. Feasibility yields nonnegative regret. An observed negative upper regret endpoint would indicate failure of the confidence event or of an implementation assumption, not a negative optimality gap.
\end{proof}

Unlike $G_d+D_{d,\epsilon}$, the new endpoint depends on the deployed policy through its paired outcomes and sample variance. It can distinguish seeds and unsuccessful policies in the same architectural tube. When $L_\phi>0$, the post-training guarantee is tighter than the analytical schedule's guarantee. This is a statement about the verified policy, not a deterministic promise that every training run improves welfare. The concentration inequality, the identity that transfers a lower improvement bound into a regret bound, and conservative policy improvement are established ideas \citep{kakade2002,pirotta2013,scherrer2015}. The contribution here is their operational connection to the continuous capital economy through \eqref{eq:r11gain} and the transfer bound \eqref{eq:r11bias}, without a state-covering certificate or a global bound on the trained network's derivatives.
